% !TeX program = lualatex
% ================================================================
% Urdu Parallel Text Version of IndicesQuickQuestions
%
% Generated by the server using the following settings:
%
%   language            Urdu (urd)
%   text direction      right to left
%   font                Noto Nastaliq Urdu
%   fallback font       Noto Serif
%   built from          latex/quickQuestions/indicesQuickQuestions.tex
%   vocabulary key      latex/quickQuestions/indicesQuickQuestions/L2/urd/indicesQuickQuestions_urduKey.tex
%   tier 3 vocabulary   green   RGB 0 128 0
%   English text        black   RGB 0 0 0
%   Urdu text           purple  RGB 112 48 160
%   layout macros       version 3
%
% Please don't edit any information in this comment block - the server
% compares it against the current settings and rebuilds this file when
% they differ, so changes here would be overwritten. However, feel free
% to make updates to the LaTeX below if you want to edit it.
% ================================================================

\documentclass[aspectratio=169]{beamer}
% ================================================================
% Copied from the English sheet this was translated from, so that
% anything it sets up for itself still works here
% ================================================================
\usepackage{amsmath}
\usepackage{amssymb}
\usepackage{tikz}
\usetikzlibrary{positioning}

% ================================================================
% Quick Questions deck: index rules.
%
% Each topic opens with five true/false statements and then runs ten
% quick questions that get gradually harder. Every question gets two
% slides: the question on its own for the mini whiteboards, then the
% same question again with the solution underneath in red.
%
% The three green topics work on the index laws themselves, first with
% positive indices, then with negative ones, then applied to algebraic
% expressions. The four amber topics move on to roots: estimating them,
% then indices of the form 1/a and a/b, and finally negative
% fractional indices.
%
% Where the working can be drawn it is: a halving ladder for the
% negative indices, number lines for the estimating, area and volume
% pictures for the square and cube roots, and flow charts for the
% fractional indices, which show the root, the power and the
% reciprocal as separate steps done in order.
% ================================================================

\setbeamertemplate{navigation symbols}{}

% A link in the bottom left of every slide back to the contents, so a
% deck this long can be jumped around in during a lesson.
\setbeamertemplate{footline}{%
  \hbox to\paperwidth{%
    \hspace{1em}%
    \hyperlink{contents}{%
      \color{gray}\scriptsize$\leftarrow$~Back to contents}%
    \hfil}%
  \vskip4pt%
}

% Topic divider.
\newcommand{\qtopic}[1]{%
  \section{#1}%
  \begin{frame}
    \vfill
    \begin{center}
      {\usebeamercolor[fg]{structure}\Huge\bfseries #1}
    \end{center}
    \vfill
  \end{frame}%
}

% \qq[<question size>][<solution size>]{<question>}{<solution>}
% The two optional sizes drop the type for the wordier questions and
% the longer solutions, so nothing runs off the slide.
\NewDocumentCommand{\qq}{O{\Huge}O{\LARGE}+m+m}{%
  \begin{frame}
    \vfill
    \begin{center}
      \begin{minipage}{0.92\textwidth}
        \centering #1 #3
      \end{minipage}
    \end{center}
    \vfill
  \end{frame}%
  \begin{frame}
    \vfill
    \begin{center}
      \begin{minipage}{0.92\textwidth}
        \centering #1 #3
        \par\vspace{0.9em}
        {\color{red}#2 #4}
      \end{minipage}
    \end{center}
    \vfill
  \end{frame}%
}

% \tf[<statement size>][<answer size>]{<statement>}{<answer>}
% As \qq, but with a standing "True or false?" heading above the
% statement. The statements hunt for the usual misconceptions, so the
% answer always says why, not just which.
\NewDocumentCommand{\tf}{O{\LARGE}O{\Large}+m+m}{%
  \begin{frame}
    \vfill
    \begin{center}
      \begin{minipage}{0.92\textwidth}
        \centering
        {\usebeamercolor[fg]{structure}\Large\bfseries True or false?}
        \par\vspace{1em}
        #1 #3
      \end{minipage}
    \end{center}
    \vfill
  \end{frame}%
  \begin{frame}
    \vfill
    \begin{center}
      \begin{minipage}{0.92\textwidth}
        \centering
        {\usebeamercolor[fg]{structure}\Large\bfseries True or false?}
        \par\vspace{1em}
        #1 #3
        \par\vspace{0.9em}
        {\color{red}#2 #4}
      \end{minipage}
    \end{center}
    \vfill
  \end{frame}%
}

% Chained working, one equation per row, so that a long run of index
% laws does not have to be squeezed onto a single line.
\newcommand{\work}[1]{%
  \begin{tabular}{@{}r@{\;=\;}l@{}}#1\end{tabular}%
}

% \flow{<start>}{<first step>}{<middle>}{<second step>}{<end>}
% A flow chart of three boxes joined by two labelled arrows. Used for
% the fractional indices, where the point is that the root and the
% power are separate steps, done one after the other.
\newcommand{\flow}[5]{%
  \begin{tikzpicture}[
      fbox/.style={draw=red!60!black, rounded corners=2pt, fill=red!8,
                   inner sep=5pt, text=red, font=\large},
      farr/.style={->, red, thick, shorten >=1pt, shorten <=1pt}]
    \node[fbox] (a) {$#1$};
    \node[fbox, right=2.4cm of a] (b) {$#3$};
    \node[fbox, right=2.4cm of b] (c) {$#5$};
    \draw[farr] (a) -- node[above,text=red,font=\small]{#2} (b);
    \draw[farr] (b) -- node[above,text=red,font=\small]{#4} (c);
  \end{tikzpicture}%
}

% \flowfour{<start>}{<step>}{<second>}{<step>}{<third>}{<step>}{<end>}
% The same chart with a third arrow, for the negative fractional
% indices, where taking the reciprocal is a step of its own.
\newcommand{\flowfour}[7]{%
  \begin{tikzpicture}[
      fbox/.style={draw=red!60!black, rounded corners=2pt, fill=red!8,
                   inner sep=4pt, text=red, font=\large},
      farr/.style={->, red, thick, shorten >=1pt, shorten <=1pt}]
    \node[fbox] (a) {$#1$};
    \node[fbox, right=2.0cm of a] (b) {$#3$};
    \node[fbox, right=2.0cm of b] (c) {$#5$};
    \node[fbox, right=2.0cm of c] (d) {$#7$};
    \draw[farr] (a) -- node[above,text=red,font=\small]{#2} (b);
    \draw[farr] (b) -- node[above,text=red,font=\small]{#4} (c);
    \draw[farr] (c) -- node[above,text=red,font=\small]{#6} (d);
  \end{tikzpicture}%
}

% The powers of 2 running down from 2^3 to 2^{-3}, halving at every
% step. The negative indices are not a new rule: they are the same
% pattern carried on past 2^0=1.
\newcommand{\halvingladder}{%
  \begin{tikzpicture}[x=1cm,y=1cm,color=red,font=\small]
    \node[anchor=east] at (0,0)    {$2^{3}$};  \node[anchor=west] at (0,0)    {$=8$};
    \node[anchor=east] at (0,-0.56) {$2^{2}$};  \node[anchor=west] at (0,-0.56) {$=4$};
    \node[anchor=east] at (0,-1.12) {$2^{1}$};  \node[anchor=west] at (0,-1.12) {$=2$};
    \node[anchor=east] at (0,-1.68) {$2^{0}$};  \node[anchor=west] at (0,-1.68) {$=1$};
    \node[anchor=east] at (0,-2.24) {$2^{-1}$}; \node[anchor=west] at (0,-2.24) {$=\tfrac{1}{2}$};
    \node[anchor=east] at (0,-2.80) {$2^{-2}$}; \node[anchor=west] at (0,-2.80) {$=\tfrac{1}{4}$};
    \node[anchor=east] at (0,-3.36) {$2^{-3}$}; \node[anchor=west] at (0,-3.36) {$=\tfrac{1}{8}$};
    \foreach \i in {0,...,5}{%
      \draw[->,red,thick] (1.2,{-0.56*\i-0.16}) -- (1.2,{-0.56*\i-0.40});}
    \node[anchor=west,font=\small] at (1.4,-1.68) {halve each time};
  \end{tikzpicture}%
}

% \rootline{<left label>}{<right label>}{<root being estimated>}{<how far along, 0 to 1>}
% A number line running between the two whole numbers either side of a
% root, with the root marked in its true position between them.
\newcommand{\rootline}[4]{%
  \begin{tikzpicture}[x=1cm,y=1cm]
    \draw[red!55!black,thick] (0,0) -- (9,0);
    \draw[red!55!black,thick] (0,-0.16) -- (0,0.16);
    \draw[red!55!black,thick] (9,-0.16) -- (9,0.16);
    \node[red!55!black,below=6pt,font=\small] at (0,0) {$#1$};
    \node[red!55!black,below=6pt,font=\small] at (9,0) {$#2$};
    \fill[red] ({9*#4},0) circle (2.6pt);
    \node[red,above=5pt,font=\small] at ({9*#4},0) {$#3$};
  \end{tikzpicture}%
}

% \areasquare{<side>}{<area>} — a square labelled with its area and the
% side length that gives it. The square root of a number is the side of
% the square with that area, which is what the picture is there to say.
\newcommand{\areasquare}[2]{%
  \begin{tikzpicture}[x=1cm,y=1cm]
    \draw[red!60!black,thick,fill=red!8] (0,0) rectangle (2.4,2.4);
    \node[red,font=\small] at (1.2,1.2) {area $#2$};
    \node[red,font=\small,below=2pt] at (1.2,0) {$#1$};
    \node[red,font=\small,left=2pt] at (0,1.2) {$#1$};
  \end{tikzpicture}%
}

% \volumecube{<edge>}{<volume>} — the matching picture for cube roots:
% the cube root of a number is the edge of the cube with that volume.
\newcommand{\volumecube}[2]{%
  \begin{tikzpicture}[x=1cm,y=1cm]
    \draw[red!60!black,thick,fill=red!8] (0,0) rectangle (2,2);
    \draw[red!60!black,thick] (0,2) -- (0.7,2.7) -- (2.7,2.7) -- (2,2);
    \draw[red!60!black,thick] (2.7,2.7) -- (2.7,0.7) -- (2,0);
    \node[red,font=\small] at (1,1) {vol. $#2$};
    \node[red,font=\small,below=2pt] at (1,0) {$#1$};
  \end{tikzpicture}%
}

% Contents-slide bullets, colour-coded by difficulty band: green for
% the three index-law topics, orange for the four root topics. The
% green is darkened because a pure green washes out on a projector.
\setbeamertemplate{section in toc}{%
  \ifnum\inserttocsectionnumber<4
    \textcolor{green!55!black}{$\bullet$}%
  \else
    \textcolor{orange}{$\bullet$}%
  \fi
  ~\inserttocsection%
}

\title{Index Rules: Quick Questions}
\subtitle{Mini whiteboard questions, with solutions}
\author{}
\date{}

% Characters this language's font has not got are borrowed from another,
% rather than being silently dropped - see docs/translated-sheet-latex.md
\directlua{%
  if luaotfload and luaotfload.add_fallback then
    luaotfload.add_fallback("eallatinfallback",
      { "Noto Serif:mode=harf;" })
  else
    tex.error("this luaotfload is too old to register a fallback font")
  end
}
\usepackage[english,bidi=basic]{babel}
\babelprovide[import]{urdu}
\babelfont[urdu]{rm}[Renderer=Harfbuzz, BoldFont={Noto Nastaliq Urdu}, ItalicFont={Noto Nastaliq Urdu}, BoldItalicFont={Noto Nastaliq Urdu}, RawFeature={fallback=eallatinfallback}]{Noto Nastaliq Urdu}
\babelfont[urdu]{sf}[Renderer=Harfbuzz, BoldFont={Noto Nastaliq Urdu}, ItalicFont={Noto Nastaliq Urdu}, BoldItalicFont={Noto Nastaliq Urdu}, RawFeature={fallback=eallatinfallback}]{Noto Nastaliq Urdu}
\newcommand{\ealtext}[1]{\foreignlanguage{urdu}{#1}}
\newcommand{\ealtextblock}[1]{{\raggedleft\foreignlanguage{urdu}{#1}\par}}
% ================================================================
% EAL layout helpers
% ================================================================
\usepackage{xcolor}

\definecolor{ealkeycolour}{RGB}{0,128,0}
\definecolor{ealenglish}{RGB}{0,0,0}
\definecolor{ealtwocolour}{RGB}{112,48,160}

\newsavebox{\ealboxen}
\newsavebox{\ealboxtr}
\newlength{\ealglosswd}
\newlength{\ealglossraise}
\setlength{\ealglossraise}{1.15em}

% a tier 3 word, in English and in translation alike
\newcommand{\ealkey}[1]{\textcolor{ealkeycolour}{#1}}
\newcommand{\ealkeytr}[1]{\textcolor{ealkeycolour}{\ealtext{#1}}}

% One translated word above one English word. Built from kernel
% primitives, never a tabular, which would break wherever the sheet
% loads array, booktabs or tabularx. The box is as wide as the wider of
% the two, so a long translation pushes its neighbours aside instead of
% overlapping them, and it claims its full height so a table row or a
% TikZ node makes room for it.
\newcommand{\ealgloss}[2]{%
  \sbox{\ealboxen}{\ealkey{#1}}%
  \sbox{\ealboxtr}{\small\textcolor{ealtwocolour}{\ealtext{#2}}}%
  \setlength{\ealglosswd}{\wd\ealboxen}%
  \ifdim\wd\ealboxtr>\ealglosswd\setlength{\ealglosswd}{\wd\ealboxtr}\fi
  \leavevmode
  \makebox[\ealglosswd][c]{%
    \makebox[0pt][c]{\usebox{\ealboxen}}%
    \makebox[0pt][c]{\raisebox{\ealglossraise}{\usebox{\ealboxtr}}}%
  }%
}

% opens up line spacing around a block of glosses, so the raised
% translations sit clear of the line above
\newenvironment{ealglossed}{\par\linespread{2.1}\selectfont\ignorespaces}{\par}

% a whole translated sentence above its English counterpart. block level,
% so it does not belong inside a tabular cell
\newcommand{\ealpara}[2]{%
  \par\addvspace{0.5\baselineskip}%
  {\color{ealtwocolour}\ealtextblock{#1}}%
  \penalty200
  {\color{ealenglish}#2\par}%
  \addvspace{0.5\baselineskip}%
}

\begin{document}

\frame{\titlepage}

\begin{frame}[label=contents]{\ealpara{فہرست}{Contents}}
  \small
  \tableofcontents
\end{frame}

%================================================================
\qtopic{\ealpara{مثبت \ealkeytr{اس} کے ساتھ \ealkeytr{اس} کے قواعد}{\ealkey{Index} rules with positive \ealkey{indices}}}
%================================================================

\tf{\ealpara{$4^{3}$ کا مطلب ہے $4\times 4\times 4$}{$4^{3}$ means $4\times 4\times 4$}}
  {\ealpara{\textbf{درست۔} \ealkeytr{اس} سے ظاہر ہوتا ہے کہ کتنے $4$ آپس میں ضرب دیے گئے ہیں، اس لیے $4^{3}=64$۔}{\textbf{True.} The \ealkey{index} counts how many $4$s are multiplied together, so $4^{3}=64$.}}

\tf{\ealpara{$2^{5}\times 2^{3}=2^{8}$}{$2^{5}\times 2^{3}=2^{8}$}}
  {\ealpara{\textbf{درست۔} پانچ $2$ کو تین $2$ سے ضرب دینے سے آٹھ $2$ حاصل ہوتے ہیں، اس لیے \ealkeytr{اس} جمع ہو جاتے ہیں۔}{\textbf{True.} Five $2$s multiplied by three $2$s gives eight $2$s, so the \ealkey{indices} are added.}}

\tf{\ealpara{$2^{5}\times 2^{3}=4^{8}$}{$2^{5}\times 2^{3}=4^{8}$}}
  {\ealpara{\textbf{غلط۔} صرف \ealkeytr{اس} جمع ہوتے ہیں؛ \ealkeytr{قاعدہ} $2$ ہی رہتا ہے۔\\[0.3em]
   جواب $2^{8}=256$ ہے، نہ کہ $4^{8}=65\,536$۔}{\textbf{False.} Only the \ealkey{indices} are added; the \ealkey{base} stays as $2$.\\[0.3em]
   The answer is $2^{8}=256$, not $4^{8}=65\,536$.}}

\tf{\ealpara{$\dfrac{7^{8}}{7^{2}}=7^{4}$}{$\dfrac{7^{8}}{7^{2}}=7^{4}$}}
  {\ealpara{\textbf{غلط۔} \ealkeytr{اس} تفریق ہوتے ہیں، تقسیم نہیں۔\\[0.3em]
   $7^{8}\div 7^{2}=7^{8-2}=7^{6}$}{\textbf{False.} The \ealkey{indices} are subtracted, not divided.\\[0.3em]
   $7^{8}\div 7^{2}=7^{8-2}=7^{6}$}}

\tf{\ealpara{$\left(5^{2}\right)^{3}=5^{5}$}{$\left(5^{2}\right)^{3}=5^{5}$}}
  {\ealpara{\textbf{غلط۔} $\left(5^{2}\right)^{3}=5^{2}\times 5^{2}\times 5^{2}$، جو چھ $5$ ہیں۔\\[0.3em]
   \ealkeytr{اس} ضرب ہوتے ہیں: $5^{6}$۔}{\textbf{False.} $\left(5^{2}\right)^{3}=5^{2}\times 5^{2}\times 5^{2}$, which is six $5$s.\\[0.3em]
   The \ealkey{indices} are multiplied: $5^{6}$.}}

\qq[\Huge][\large]{\ealpara{\ealkeytr{سادہ کرنا} $2^{3}\times 2^{4}$}{\ealkey{Simplify} $2^{3}\times 2^{4}$}}
  {\ealpara{$(2\times 2\times 2)\times(2\times 2\times 2\times 2)$ --- سات $2$\\[0.4em]
   $2^{3}\times 2^{4}=2^{3+4}=2^{7}$}{$(2\times 2\times 2)\times(2\times 2\times 2\times 2)$ --- seven $2$s\\[0.4em]
   $2^{3}\times 2^{4}=2^{3+4}=2^{7}$}}

\qq{\ealpara{\ealkeytr{سادہ کرنا} $5^{6}\div 5^{2}$}{\ealkey{Simplify} $5^{6}\div 5^{2}$}}
  {\ealpara{$5^{6-2}=5^{4}$}{$5^{6-2}=5^{4}$}}

\qq{\ealpara{\ealkeytr{سادہ کرنا} $7^{5}\times 7$}{\ealkey{Simplify} $7^{5}\times 7$}}
  {\ealpara{$7$ اکیلا $7^{1}$ ہے۔\\[0.3em]
   $7^{5+1}=7^{6}$}{$7$ on its own is $7^{1}$.\\[0.3em]
   $7^{5+1}=7^{6}$}}

\qq{\ealpara{\ealkeytr{سادہ کرنا} $\left(3^{2}\right)^{4}$}{\ealkey{Simplify} $\left(3^{2}\right)^{4}$}}
  {\ealpara{$3^{2\times 4}=3^{8}$}{$3^{2\times 4}=3^{8}$}}

\qq{\ealpara{\ealkeytr{سادہ کرنا} $4^{3}\times 4^{2}\times 4^{5}$}{\ealkey{Simplify} $4^{3}\times 4^{2}\times 4^{5}$}}
  {\ealpara{$4^{3+2+5}=4^{10}$}{$4^{3+2+5}=4^{10}$}}

\qq[\Huge][\LARGE]{\ealpara{\ealkeytr{سادہ کرنا} $\dfrac{6^{8}}{6^{3}\times 6^{2}}$}{\ealkey{Simplify} $\dfrac{6^{8}}{6^{3}\times 6^{2}}$}}
  {\work{$6^{3}\times 6^{2}$ & $6^{5}$\\[0.3em]
         $\dfrac{6^{8}}{6^{5}}$ & $6^{8-5}=6^{3}$}}

\qq[\Huge][\LARGE]{\ealpara{\ealkeytr{سادہ کرنا} $\left(2^{3}\right)^{2}\times 2^{4}$}{\ealkey{Simplify} $\left(2^{3}\right)^{2}\times 2^{4}$}}
  {\work{$\left(2^{3}\right)^{2}$ & $2^{6}$\\[0.3em]
         $2^{6}\times 2^{4}$ & $2^{10}$}}

\qq[\Huge][\LARGE]{\ealpara{\ealkeytr{سادہ کرنا} $\dfrac{\left(5^{3}\right)^{4}}{5^{10}}$}{\ealkey{Simplify} $\dfrac{\left(5^{3}\right)^{4}}{5^{10}}$}}
  {\work{$\left(5^{3}\right)^{4}$ & $5^{12}$\\[0.3em]
         $\dfrac{5^{12}}{5^{10}}$ & $5^{12-10}=5^{2}$}}

\qq[\LARGE][\LARGE]{\ealpara{$3^{n}\times 3^{5}=3^{12}$\\[0.4em]$n$ کی قدر معلوم کریں}{$3^{n}\times 3^{5}=3^{12}$\\[0.4em]Find the value of $n$}}
  {\ealpara{\ealkeytr{اس} جمع ہوتے ہیں، اس لیے $n+5=12$۔\\[0.3em]
   $n=7$}{The \ealkey{indices} add, so $n+5=12$.\\[0.3em]
   $n=7$}}

\qq[\LARGE][\LARGE]{\ealpara{$8^{4}$ کو $2$ کی \ealkeytr{قوت} کے طور پر لکھیں}{Write $8^{4}$ as a \ealkey{power} of $2$}}
  {\ealpara{$8=2^{3}$، اس لیے $8^{4}=\left(2^{3}\right)^{4}$۔\\[0.3em]
   $8^{4}=2^{12}$}{$8=2^{3}$, so $8^{4}=\left(2^{3}\right)^{4}$.\\[0.3em]
   $8^{4}=2^{12}$}}

%================================================================
\qtopic{\ealpara{منفی \ealkeytr{اس} کے ساتھ \ealkeytr{اس} کے قواعد}{\ealkey{Index} rules with negative \ealkey{indices}}}
%================================================================

\tf[\Large][\normalsize]{\ealpara{$2^{-3}$ ایک منفی عدد ہے}{$2^{-3}$ is a negative number}}
  {\ealpara{\textbf{غلط۔} منفی \ealkeytr{اس} \ealkeytr{مقلوب} دیتا ہے، منفی جواب نہیں۔\\[0.3em]
   $2^{-3}=\dfrac{1}{2^{3}}=\dfrac{1}{8}$، جو مثبت ہے۔
   \par\vspace{0.4em}
   \halvingladder}{\textbf{False.} A negative \ealkey{index} gives a \ealkey{reciprocal}, not a negative answer.\\[0.3em]
   $2^{-3}=\dfrac{1}{2^{3}}=\dfrac{1}{8}$, which is positive.
   \par\vspace{0.4em}
   \halvingladder}}

\tf{\ealpara{$5^{-1}=\dfrac{1}{5}$}{$5^{-1}=\dfrac{1}{5}$}}
  {\ealpara{\textbf{درست۔} $-1$ کا \ealkeytr{اس} \ealkeytr{قاعدہ} کا \ealkeytr{مقلوب} دیتا ہے۔}{\textbf{True.} An \ealkey{index} of $-1$ gives the \ealkey{reciprocal} of the \ealkey{base}.}}

\tf{\ealpara{$4^{-2}=\dfrac{1}{8}$}{$4^{-2}=\dfrac{1}{8}$}}
  {\ealpara{\textbf{غلط۔} $-2$ ایک \ealkeytr{اس} ہے، ضارب نہیں۔\\[0.3em]
   $4^{-2}=\dfrac{1}{4^{2}}=\dfrac{1}{16}$}{\textbf{False.} The $-2$ is an \ealkey{index}, not a multiplier.\\[0.3em]
   $4^{-2}=\dfrac{1}{4^{2}}=\dfrac{1}{16}$}}

\tf{\ealpara{$2^{4}\div 2^{7}=2^{-3}$}{$2^{4}\div 2^{7}=2^{-3}$}}
  {\ealpara{\textbf{درست۔} \ealkeytr{اس} تفریق ہوتے ہیں: $4-7=-3$۔}{\textbf{True.} The \ealkey{indices} are subtracted: $4-7=-3$.}}

\tf{\ealpara{$2^{-3}\times 2^{5}=2^{-15}$}{$2^{-3}\times 2^{5}=2^{-15}$}}
  {\ealpara{\textbf{غلط۔} \ealkeytr{اس} جمع ہوتے ہیں، ضرب نہیں۔\\[0.3em]
   $2^{-3+5}=2^{2}=4$}{\textbf{False.} The \ealkey{indices} are added, not multiplied.\\[0.3em]
   $2^{-3+5}=2^{2}=4$}}

\qq[\Huge][\large]{\ealpara{$2^{-1}$ کو \ealkeytr{کسر} کی شکل میں لکھیں}{Write $2^{-1}$ as a \ealkey{fraction}}}
  {\ealpara{$2^{-1}=\dfrac{1}{2}$
   \par\vspace{0.5em}
   \halvingladder}{$2^{-1}=\dfrac{1}{2}$
   \par\vspace{0.5em}
   \halvingladder}}

\qq{\ealpara{$3^{-2}$ کو \ealkeytr{کسر} کی شکل میں لکھیں}{Write $3^{-2}$ as a \ealkey{fraction}}}
  {\ealpara{$3^{-2}=\dfrac{1}{3^{2}}=\dfrac{1}{9}$}{$3^{-2}=\dfrac{1}{3^{2}}=\dfrac{1}{9}$}}

\qq[\Huge][\LARGE]{\ealpara{$\dfrac{1}{5^{3}}$ کو منفی \ealkeytr{اس} کا استعمال کرتے ہوئے لکھیں}{Write $\dfrac{1}{5^{3}}$ using a negative \ealkey{index}}}
  {\ealpara{$\dfrac{1}{5^{3}}=5^{-3}$}{$\dfrac{1}{5^{3}}=5^{-3}$}}

\qq[\LARGE][\LARGE]{\ealpara{$10^{-3}$ کو \ealkeytr{اعشاریہ} کی شکل میں لکھیں}{Write $10^{-3}$ as a \ealkey{decimal}}}
  {\ealpara{$10^{-3}=\dfrac{1}{10^{3}}=\dfrac{1}{1000}=0.001$}{$10^{-3}=\dfrac{1}{10^{3}}=\dfrac{1}{1000}=0.001$}}

\qq{\ealpara{\ealkeytr{سادہ کرنا} $2^{3}\times 2^{-5}$}{\ealkey{Simplify} $2^{3}\times 2^{-5}$}}
  {\ealpara{$2^{3+(-5)}=2^{-2}$}{$2^{3+(-5)}=2^{-2}$}}

\qq{\ealpara{\ealkeytr{سادہ کرنا} $7^{-2}\div 7^{3}$}{\ealkey{Simplify} $7^{-2}\div 7^{3}$}}
  {\ealpara{$7^{-2-3}=7^{-5}$}{$7^{-2-3}=7^{-5}$}}

\qq{\ealpara{\ealkeytr{سادہ کرنا} $\left(4^{-2}\right)^{3}$}{\ealkey{Simplify} $\left(4^{-2}\right)^{3}$}}
  {\ealpara{$4^{-2\times 3}=4^{-6}$}{$4^{-2\times 3}=4^{-6}$}}

\qq[\Huge][\LARGE]{\ealpara{$\dfrac{5^{-3}}{5^{-7}}$ کا حساب لگائیں}{Work out $\dfrac{5^{-3}}{5^{-7}}$}}
  {\work{$5^{-3-(-7)}$ & $5^{-3+7}$\\[0.3em]
         $5^{4}$ & $625$}}

\qq[\Huge][\LARGE]{\ealpara{$\left(\dfrac{2}{3}\right)^{-2}$ کا حساب لگائیں}{Work out $\left(\dfrac{2}{3}\right)^{-2}$}}
  {\ealpara{منفی \ealkeytr{اس} \ealkeytr{کسر} کا \ealkeytr{مقلوب} دیتا ہے۔\\[0.4em]
   $\left(\dfrac{3}{2}\right)^{2}=\dfrac{9}{4}$}{A negative \ealkey{index} gives the \ealkey{reciprocal} of the \ealkey{fraction}.\\[0.4em]
   $\left(\dfrac{3}{2}\right)^{2}=\dfrac{9}{4}$}}

\qq[\LARGE][\LARGE]{\ealpara{$2^{n}\times 2^{-5}=2^{-1}$\\[0.4em]$n$ کی قدر معلوم کریں}{$2^{n}\times 2^{-5}=2^{-1}$\\[0.4em]Find the value of $n$}}
  {\ealpara{\ealkeytr{اس} جمع ہوتے ہیں، اس لیے $n-5=-1$۔\\[0.3em]
   $n=4$}{The \ealkey{indices} add, so $n-5=-1$.\\[0.3em]
   $n=4$}}

%================================================================
\qtopic{\ealpara{\ealkeytr{قوت نما کا قانون} کے ساتھ \ealkeytr{اظہار} کو \ealkeytr{سادہ کرنا}}{\ealkey{Simplifying} \ealkey{expressions} with \ealkey{index law}s}}
%================================================================

\tf{\ealpara{$x^{3}\times x^{4}=x^{12}$}{$x^{3}\times x^{4}=x^{12}$}}
  {\ealpara{\textbf{غلط۔} \ealkeytr{اس} جمع ہوتے ہیں، ضرب نہیں۔\\[0.3em]
   $x^{3}\times x^{4}=x^{7}$}{\textbf{False.} The \ealkey{indices} are added, not multiplied.\\[0.3em]
   $x^{3}\times x^{4}=x^{7}$}}

\tf{\ealpara{$2a^{3}\times 3a^{2}=6a^{5}$}{$2a^{3}\times 3a^{2}=6a^{5}$}}
  {\ealpara{\textbf{درست۔} اعداد ضرب ہوتے ہیں اور \ealkeytr{اس} جمع ہوتے ہیں۔}{\textbf{True.} The numbers are multiplied and the \ealkey{indices} are added.}}

\tf{\ealpara{$\left(2x^{3}\right)^{2}=2x^{6}$}{$\left(2x^{3}\right)^{2}=2x^{6}$}}
  {\ealpara{\textbf{غلط۔} بریکٹ کے اندر ہر چیز کا \ealkeytr{مربع} لیا جاتا ہے، بشمول $2$۔\\[0.3em]
   $\left(2x^{3}\right)^{2}=2^{2}\times x^{6}=4x^{6}$}{\textbf{False.} Everything inside the bracket is \ealkey{squared}, including the $2$.\\[0.3em]
   $\left(2x^{3}\right)^{2}=2^{2}\times x^{6}=4x^{6}$}}

\tf{\ealpara{$x^{5}+x^{3}=x^{8}$}{$x^{5}+x^{3}=x^{8}$}}
  {\ealpara{\textbf{غلط۔} \ealkeytr{اس} جمع ہوتے ہیں جب \ealkeytr{رکن} ضرب دیے جاتے ہیں، جب وہ جمع کیے جاتے ہیں تو نہیں۔\\[0.3em]
   $x^{5}+x^{3}$ کو \ealkeytr{سادہ} نہیں کیا جا سکتا۔}{\textbf{False.} \ealkey{Indices} are added when \ealkey{terms} are multiplied, not when they are added.\\[0.3em]
   $x^{5}+x^{3}$ cannot be \ealkey{simplified}.}}

\tf{\ealpara{$\dfrac{12x^{5}}{3x^{5}}=4$}{$\dfrac{12x^{5}}{3x^{5}}=4$}}
  {\ealpara{\textbf{درست۔} $12\div 3=4$ اور $x^{5}\div x^{5}=x^{0}=1$۔}{\textbf{True.} $12\div 3=4$ and $x^{5}\div x^{5}=x^{0}=1$.}}

\qq{\ealpara{\ealkeytr{سادہ کرنا} $x^{3}\times x^{4}$}{\ealkey{Simplify} $x^{3}\times x^{4}$}}
  {\ealpara{$x^{3+4}=x^{7}$}{$x^{3+4}=x^{7}$}}

\qq{\ealpara{\ealkeytr{سادہ کرنا} $y^{8}\div y^{3}$}{\ealkey{Simplify} $y^{8}\div y^{3}$}}
  {\ealpara{$y^{8-3}=y^{5}$}{$y^{8-3}=y^{5}$}}

\qq{\ealpara{\ealkeytr{سادہ کرنا} $3a^{2}\times 4a^{5}$}{\ealkey{Simplify} $3a^{2}\times 4a^{5}$}}
  {\ealpara{$3\times 4=12$ اور $a^{2+5}=a^{7}$\\[0.3em]
   $12a^{7}$}{$3\times 4=12$ and $a^{2+5}=a^{7}$\\[0.3em]
   $12a^{7}$}}

\qq[\Huge][\LARGE]{\ealpara{\ealkeytr{سادہ کرنا} $\dfrac{20b^{7}}{5b^{2}}$}{\ealkey{Simplify} $\dfrac{20b^{7}}{5b^{2}}$}}
  {\ealpara{$20\div 5=4$ اور $b^{7-2}=b^{5}$\\[0.3em]
   $4b^{5}$}{$20\div 5=4$ and $b^{7-2}=b^{5}$\\[0.3em]
   $4b^{5}$}}

\qq{\ealpara{\ealkeytr{سادہ کرنا} $\left(x^{4}\right)^{3}$}{\ealkey{Simplify} $\left(x^{4}\right)^{3}$}}
  {\ealpara{$x^{4\times 3}=x^{12}$}{$x^{4\times 3}=x^{12}$}}

\qq[\Huge][\LARGE]{\ealpara{\ealkeytr{سادہ کرنا} $\left(3m^{2}\right)^{3}$}{\ealkey{Simplify} $\left(3m^{2}\right)^{3}$}}
  {\ealpara{$3$ کا بھی \ealkeytr{مکعب} لیا جاتا ہے اور $m^{2}$ کا بھی۔\\[0.3em]
   $3^{3}\times m^{6}=27m^{6}$}{The $3$ is \ealkey{cubed} as well as the $m^{2}$.\\[0.3em]
   $3^{3}\times m^{6}=27m^{6}$}}

\qq[\LARGE][\LARGE]{\ealpara{\ealkeytr{سادہ کرنا} $5p^{3}q^{2}\times 4p^{2}q^{5}$}{\ealkey{Simplify} $5p^{3}q^{2}\times 4p^{2}q^{5}$}}
  {\ealpara{$5\times 4=20$، $p^{3+2}=p^{5}$، $q^{2+5}=q^{7}$\\[0.3em]
   $20p^{5}q^{7}$}{$5\times 4=20$, $p^{3+2}=p^{5}$, $q^{2+5}=q^{7}$\\[0.3em]
   $20p^{5}q^{7}$}}

\qq[\LARGE][\LARGE]{\ealpara{\ealkeytr{سادہ کرنا} $\dfrac{18x^{6}y^{4}}{6x^{2}y^{4}}$}{\ealkey{Simplify} $\dfrac{18x^{6}y^{4}}{6x^{2}y^{4}}$}}
  {\ealpara{$18\div 6=3$، $x^{6-2}=x^{4}$، $y^{4-4}=y^{0}=1$\\[0.3em]
   $3x^{4}$}{$18\div 6=3$, $x^{6-2}=x^{4}$, $y^{4-4}=y^{0}=1$\\[0.3em]
   $3x^{4}$}}

\qq[\LARGE][\LARGE]{\ealpara{\ealkeytr{سادہ کرنا} $\dfrac{\left(2x^{3}\right)^{4}}{8x^{5}}$}{\ealkey{Simplify} $\dfrac{\left(2x^{3}\right)^{4}}{8x^{5}}$}}
  {\work{$\left(2x^{3}\right)^{4}$ & $16x^{12}$\\[0.3em]
         $\dfrac{16x^{12}}{8x^{5}}$ & $2x^{7}$}}

\qq[\Large][\LARGE]{\ealpara{\ealkeytr{سادہ کرنا} $\dfrac{6a^{5}b^{2}\times 3a^{2}b^{4}}{9a^{4}b^{3}}$}{\ealkey{Simplify} $\dfrac{6a^{5}b^{2}\times 3a^{2}b^{4}}{9a^{4}b^{3}}$}}
  {\work{$6a^{5}b^{2}\times 3a^{2}b^{4}$ & $18a^{7}b^{6}$\\[0.3em]
         $\dfrac{18a^{7}b^{6}}{9a^{4}b^{3}}$ & $2a^{3}b^{3}$}}

%================================================================
\qtopic{\ealpara{\ealkeytr{جذر} اور \ealkeytr{قوت} کا \ealkeytr{تخمینہ}}{\ealkey{Estimating} \ealkey{roots} and \ealkey{powers}}}
%================================================================

\tf[\LARGE][\large]{\ealpara{$\sqrt{50}$ $7$ اور $8$ کے درمیان ہے}{$\sqrt{50}$ lies between $7$ and $8$}}
  {\ealpara{\textbf{درست۔} $7^{2}=49$ اور $8^{2}=64$، اور $50$ ان دو \ealkeytr{مربع} اعداد کے درمیان ہے۔
   \par\vspace{0.6em}
   \rootline{\sqrt{49}=7}{\sqrt{64}=8}{\sqrt{50}}{0.071}}{\textbf{True.} $7^{2}=49$ and $8^{2}=64$, and $50$ lies between those two \ealkey{square} numbers.
   \par\vspace{0.6em}
   \rootline{\sqrt{49}=7}{\sqrt{64}=8}{\sqrt{50}}{0.071}}}

\tf{\ealpara{$\sqrt{20}=10$}{$\sqrt{20}=10$}}
  {\ealpara{\textbf{غلط۔} \ealkeytr{جذر} لینا آدھا کرنا نہیں ہے۔\\[0.3em]
   $10^{2}=100$، $20$ نہیں۔ دراصل $\sqrt{20}$ $4$ اور $5$ کے درمیان ہے۔}{\textbf{False.} \ealkey{Square} rooting is not halving.\\[0.3em]
   $10^{2}=100$, not $20$. In fact $\sqrt{20}$ is between $4$ and $5$.}}

\tf{\ealpara{$\sqrt{1000}$ تقریباً $100$ ہے}{$\sqrt{1000}$ is about $100$}}
  {\ealpara{\textbf{غلط۔} $100^{2}=10\,000$، $1000$ نہیں۔\\[0.3em]
   $31^{2}=961$ اور $32^{2}=1024$، اس لیے $\sqrt{1000}$ تقریباً $31.6$ ہے۔}{\textbf{False.} $100^{2}=10\,000$, not $1000$.\\[0.3em]
   $31^{2}=961$ and $32^{2}=1024$, so $\sqrt{1000}$ is about $31.6$.}}

\tf[\LARGE][\large]{\ealpara{$\sqrt[3]{30}$ $3$ اور $4$ کے درمیان ہے}{$\sqrt[3]{30}$ lies between $3$ and $4$}}
  {\ealpara{\textbf{درست۔} $3^{3}=27$ اور $4^{3}=64$، اور $30$ ان دو \ealkeytr{مکعب} اعداد کے درمیان ہے۔
   \par\vspace{0.6em}
   \rootline{\sqrt[3]{27}=3}{\sqrt[3]{64}=4}{\sqrt[3]{30}}{0.107}}{\textbf{True.} $3^{3}=27$ and $4^{3}=64$, and $30$ lies between those two \ealkey{cube} numbers.
   \par\vspace{0.6em}
   \rootline{\sqrt[3]{27}=3}{\sqrt[3]{64}=4}{\sqrt[3]{30}}{0.107}}}

\tf{\ealpara{$\sqrt{0.64}$ $0.64$ سے چھوٹا ہے}{$\sqrt{0.64}$ is smaller than $0.64$}}
  {\ealpara{\textbf{غلط۔} \ealkeytr{جذر} لینا کسی عدد کو صرف اس وقت چھوٹا کرتا ہے جب وہ $1$ سے بڑا ہو۔\\[0.3em]
   $0.8^{2}=0.64$، اس لیے $\sqrt{0.64}=0.8$، جو $0.64$ سے بڑا ہے۔}{\textbf{False.} \ealkey{Square} rooting only makes a number smaller when it is bigger than $1$.\\[0.3em]
   $0.8^{2}=0.64$, so $\sqrt{0.64}=0.8$, which is larger than $0.64$.}}

\qq[\Huge][\large]{\ealpara{$\sqrt{30}$ کون سے دو \ealkeytr{مکمل عدد} کے درمیان ہے؟}{Between which two \ealkey{whole number}s\\[0.3em]does $\sqrt{30}$ lie?}}
  {\ealpara{$5$ اور $6$ کے درمیان، کیونکہ $5^{2}=25$ اور $6^{2}=36$۔
   \par\vspace{0.6em}
   \rootline{\sqrt{25}=5}{\sqrt{36}=6}{\sqrt{30}}{0.477}}{Between $5$ and $6$, because $5^{2}=25$ and $6^{2}=36$.
   \par\vspace{0.6em}
   \rootline{\sqrt{25}=5}{\sqrt{36}=6}{\sqrt{30}}{0.477}}}

\qq[\Huge][\large]{\ealpara{$\sqrt{75}$ کون سے دو \ealkeytr{مکمل عدد} کے درمیان ہے؟}{Between which two \ealkey{whole number}s\\[0.3em]does $\sqrt{75}$ lie?}}
  {\ealpara{$8$ اور $9$ کے درمیان، کیونکہ $8^{2}=64$ اور $9^{2}=81$۔
   \par\vspace{0.6em}
   \rootline{\sqrt{64}=8}{\sqrt{81}=9}{\sqrt{75}}{0.660}}{Between $8$ and $9$, because $8^{2}=64$ and $9^{2}=81$.
   \par\vspace{0.6em}
   \rootline{\sqrt{64}=8}{\sqrt{81}=9}{\sqrt{75}}{0.660}}}

\qq[\LARGE][\LARGE]{\ealpara{$3.9^{3}$ کا \ealkeytr{تخمینہ} لگائیں}{\ealkey{Estimate} $3.9^{3}$}}
  {\ealpara{$3.9$ کو $4$ تک گول کریں۔\\[0.3em]
   $4^{3}=64$، اس لیے $3.9^{3}$ تقریباً $64$ ہے۔}{Round $3.9$ to $4$.\\[0.3em]
   $4^{3}=64$, so $3.9^{3}$ is about $64$.}}

\qq[\Huge][\large]{\ealpara{$\sqrt[3]{50}$ کون سے دو \ealkeytr{مکمل عدد} کے درمیان ہے؟}{Between which two \ealkey{whole number}s\\[0.3em]does $\sqrt[3]{50}$ lie?}}
  {\ealpara{$3$ اور $4$ کے درمیان، کیونکہ $3^{3}=27$ اور $4^{3}=64$۔
   \par\vspace{0.6em}
   \rootline{\sqrt[3]{27}=3}{\sqrt[3]{64}=4}{\sqrt[3]{50}}{0.684}}{Between $3$ and $4$, because $3^{3}=27$ and $4^{3}=64$.
   \par\vspace{0.6em}
   \rootline{\sqrt[3]{27}=3}{\sqrt[3]{64}=4}{\sqrt[3]{50}}{0.684}}}

\qq[\LARGE][\large]{\ealpara{$\sqrt{17}$ کا \ealkeytr{تخمینہ} ایک \ealkeytr{اعشاریہ مقام} تک لگائیں}{\ealkey{Estimate} $\sqrt{17}$\\[0.3em]to one \ealkey{decimal place}}}
  {\ealpara{$\sqrt{16}=4$، اس لیے جواب $4$ سے ذرا اوپر ہے۔\\[0.3em]
   $4.1^{2}=16.81$ اور $4.2^{2}=17.64$، اس لیے $\sqrt{17}\approx 4.1$۔}{$\sqrt{16}=4$, so the answer is just above $4$.\\[0.3em]
   $4.1^{2}=16.81$ and $4.2^{2}=17.64$, so $\sqrt{17}\approx 4.1$.}}

\qq[\LARGE][\large]{\ealpara{$\sqrt{90}$ کا \ealkeytr{تخمینہ} ایک \ealkeytr{اعشاریہ مقام} تک لگائیں}{\ealkey{Estimate} $\sqrt{90}$\\[0.3em]to one \ealkey{decimal place}}}
  {\ealpara{$\sqrt{81}=9$ اور $\sqrt{100}=10$، اس لیے جواب $9$ اور $10$ کے درمیان ہے۔\\[0.3em]
   $9.5^{2}=90.25$، اس لیے $\sqrt{90}\approx 9.5$۔}{$\sqrt{81}=9$ and $\sqrt{100}=10$, so the answer is between $9$ and $10$.\\[0.3em]
   $9.5^{2}=90.25$, so $\sqrt{90}\approx 9.5$.}}

\qq[\LARGE][\LARGE]{\ealpara{کون بڑا ہے، $\sqrt{40}$ یا $6.5$؟}{Which is larger,\\[0.3em]$\sqrt{40}$ or $6.5$?}}
  {\ealpara{دونوں کا \ealkeytr{مربع} لیں: $\left(\sqrt{40}\right)^{2}=40$ اور $6.5^{2}=42.25$۔\\[0.3em]
   $6.5$ بڑا ہے۔}{\ealkey{Square} both: $\left(\sqrt{40}\right)^{2}=40$ and $6.5^{2}=42.25$.\\[0.3em]
   $6.5$ is larger.}}

\qq[\LARGE][\large]{\ealpara{$\sqrt{2000}$ کا \ealkeytr{تخمینہ} قریب ترین \ealkeytr{مکمل عدد} تک لگائیں}{\ealkey{Estimate} $\sqrt{2000}$\\[0.3em]to the nearest \ealkey{whole number}}}
  {\ealpara{$44^{2}=1936$ اور $45^{2}=2025$۔\\[0.3em]
   $2000$، $2025$ کے زیادہ قریب ہے، اس لیے $\sqrt{2000}\approx 45$۔}{$44^{2}=1936$ and $45^{2}=2025$.\\[0.3em]
   $2000$ is nearer to $2025$, so $\sqrt{2000}\approx 45$.}}

\qq[\LARGE][\LARGE]{\ealpara{$\dfrac{\sqrt{101}}{4.9}$ کا \ealkeytr{تخمینہ} لگائیں}{\ealkey{Estimate} $\dfrac{\sqrt{101}}{4.9}$}}
  {\ealpara{ہر حصے کو گول کریں: $\sqrt{101}\approx\sqrt{100}=10$ اور $4.9\approx 5$۔\\[0.3em]
   $\dfrac{10}{5}=2$}{Round each part: $\sqrt{101}\approx\sqrt{100}=10$ and $4.9\approx 5$.\\[0.3em]
   $\dfrac{10}{5}=2$}}

\qq[\Large][\large]{\ealpara{ایک \ealkeytr{مربع} آنگن کا \ealkeytr{رقبہ} $150\ \mathrm{m^{2}}$ ہے۔\\[0.4em]
   ایک طرف کی لمبائی کا \ealkeytr{تخمینہ} لگائیں۔}{A \ealkey{square} patio has an \ealkey{area} of $150\ \mathrm{m^{2}}$.\\[0.4em]
   \ealkey{Estimate} the length of one side.}}
  {\ealpara{طرف $\sqrt{150}$ ہے۔\\[0.3em]
   $12^{2}=144$ اور $13^{2}=169$، اس لیے طرف $12$ سے ذرا زیادہ ہے۔\\[0.3em]
   $12.2^{2}=148.84$، اس لیے طرف تقریباً $12.2\ \mathrm{m}$ ہے۔}{The side is $\sqrt{150}$.\\[0.3em]
   $12^{2}=144$ and $13^{2}=169$, so the side is just over $12$.\\[0.3em]
   $12.2^{2}=148.84$, so the side is about $12.2\ \mathrm{m}$.}}

%================================================================
\qtopic{\ealpara{$\frac{1}{a}$ کی شکل کے \ealkeytr{اس}}{\ealkey{Indices} of the form $\frac{1}{a}$}}
%================================================================

\tf[\LARGE][\large]{\ealpara{$9^{\frac{1}{2}}=3$}{$9^{\frac{1}{2}}=3$}}
  {\ealpara{\textbf{درست۔} $\frac{1}{2}$ کا \ealkeytr{اس} \ealkeytr{جذر} کی نشاندہی کرتا ہے، اور $3^{2}=9$۔
   \par\vspace{0.5em}
   \areasquare{3}{9}}{\textbf{True.} An \ealkey{index} of $\frac{1}{2}$ means the \ealkey{square root}, and $3^{2}=9$.
   \par\vspace{0.5em}
   \areasquare{3}{9}}}

\tf{\ealpara{$9^{\frac{1}{2}}=4.5$}{$9^{\frac{1}{2}}=4.5$}}
  {\ealpara{\textbf{غلط۔} \ealkeytr{اس} ضارب نہیں ہے، اس لیے یہ $9$ کا آدھا نہیں ہے۔\\[0.3em]
   $9^{\frac{1}{2}}=\sqrt{9}=3$}{\textbf{False.} The \ealkey{index} is not a multiplier, so this is not half of $9$.\\[0.3em]
   $9^{\frac{1}{2}}=\sqrt{9}=3$}}

\tf[\LARGE][\large]{\ealpara{$8^{\frac{1}{3}}=2$}{$8^{\frac{1}{3}}=2$}}
  {\ealpara{\textbf{درست۔} $\frac{1}{3}$ کا \ealkeytr{اس} \ealkeytr{مکعب جذر} کی نشاندہی کرتا ہے، اور $2^{3}=8$۔
   \par\vspace{0.5em}
   \volumecube{2}{8}}{\textbf{True.} An \ealkey{index} of $\frac{1}{3}$ means the \ealkey{cube root}, and $2^{3}=8$.
   \par\vspace{0.5em}
   \volumecube{2}{8}}}

\tf{\ealpara{$16^{\frac{1}{4}}=4$}{$16^{\frac{1}{4}}=4$}}
  {\ealpara{\textbf{غلط۔} $\frac{1}{4}$ کا \ealkeytr{اس} \ealkeytr{چوتھا جذر} کی نشاندہی کرتا ہے، \ealkeytr{جذر} کا نہیں۔\\[0.3em]
   $2^{4}=16$، اس لیے $16^{\frac{1}{4}}=2$۔ یہ $16^{\frac{1}{2}}$ ہے جو $4$ کے برابر ہے۔}{\textbf{False.} An \ealkey{index} of $\frac{1}{4}$ means the \ealkey{fourth root}, not the \ealkey{square root}.\\[0.3em]
   $2^{4}=16$, so $16^{\frac{1}{4}}=2$. It is $16^{\frac{1}{2}}$ that equals $4$.}}

\tf{\ealpara{$100^{\frac{1}{3}}=10$}{$100^{\frac{1}{3}}=10$}}
  {\ealpara{\textbf{غلط۔} $10^{3}=1000$، $100$ نہیں۔\\[0.3em]
   $\sqrt[3]{100}$ تقریباً $4.6$ ہے۔ یہ $100^{\frac{1}{2}}$ ہے جو $10$ کے برابر ہے۔}{\textbf{False.} $10^{3}=1000$, not $100$.\\[0.3em]
   $\sqrt[3]{100}$ is about $4.6$. It is $100^{\frac{1}{2}}$ that equals $10$.}}

\qq[\Huge][\large]{\ealpara{$25^{\frac{1}{2}}$ کا حساب لگائیں}{Work out $25^{\frac{1}{2}}$}}
  {\ealpara{$25^{\frac{1}{2}}=\sqrt{25}=5$
   \par\vspace{0.5em}
   \areasquare{5}{25}}{$25^{\frac{1}{2}}=\sqrt{25}=5$
   \par\vspace{0.5em}
   \areasquare{5}{25}}}

\qq{\ealpara{$64^{\frac{1}{2}}$ کا حساب لگائیں}{Work out $64^{\frac{1}{2}}$}}
  {\ealpara{$\sqrt{64}=8$، کیونکہ $8^{2}=64$۔}{$\sqrt{64}=8$, because $8^{2}=64$.}}

\qq[\Huge][\large]{\ealpara{$27^{\frac{1}{3}}$ کا حساب لگائیں}{Work out $27^{\frac{1}{3}}$}}
  {\ealpara{$\sqrt[3]{27}=3$، کیونکہ $3^{3}=27$۔
   \par\vspace{0.5em}
   \volumecube{3}{27}}{$\sqrt[3]{27}=3$, because $3^{3}=27$.
   \par\vspace{0.5em}
   \volumecube{3}{27}}}

\qq[\Huge][\LARGE]{\ealpara{$81^{\frac{1}{4}}$ کا حساب لگائیں}{Work out $81^{\frac{1}{4}}$}}
  {\ealpara{\ealkeytr{چوتھا جذر}: $3\times 3\times 3\times 3=81$۔\\[0.3em]
   $81^{\frac{1}{4}}=3$}{The \ealkey{fourth root}: $3\times 3\times 3\times 3=81$.\\[0.3em]
   $81^{\frac{1}{4}}=3$}}

\qq[\Huge][\LARGE]{\ealpara{$1000^{\frac{1}{3}}$ کا حساب لگائیں}{Work out $1000^{\frac{1}{3}}$}}
  {\ealpara{$\sqrt[3]{1000}=10$، کیونکہ $10^{3}=1000$۔}{$\sqrt[3]{1000}=10$, because $10^{3}=1000$.}}

\qq[\Huge][\LARGE]{\ealpara{$\sqrt[5]{x}$ کو \ealkeytr{اس} کا استعمال کرتے ہوئے لکھیں}{Write $\sqrt[5]{x}$ using an \ealkey{index}}}
  {\ealpara{$\sqrt[5]{x}=x^{\frac{1}{5}}$}{$\sqrt[5]{x}=x^{\frac{1}{5}}$}}

\qq[\Huge][\LARGE]{\ealpara{$\left(\dfrac{1}{16}\right)^{\frac{1}{2}}$ کا حساب لگائیں}{Work out $\left(\dfrac{1}{16}\right)^{\frac{1}{2}}$}}
  {\ealpara{\ealkeytr{صورت} اور \ealkeytr{مخرج} کا \ealkeytr{جذر} لیں۔\\[0.3em]
   $\dfrac{\sqrt{1}}{\sqrt{16}}=\dfrac{1}{4}$}{\ealkey{Square root} the \ealkey{numerator} and the \ealkey{denominator}.\\[0.3em]
   $\dfrac{\sqrt{1}}{\sqrt{16}}=\dfrac{1}{4}$}}

\qq[\Huge][\LARGE]{\ealpara{$\left(\dfrac{8}{27}\right)^{\frac{1}{3}}$ کا حساب لگائیں}{Work out $\left(\dfrac{8}{27}\right)^{\frac{1}{3}}$}}
  {\ealpara{\ealkeytr{صورت} اور \ealkeytr{مخرج} کا \ealkeytr{مکعب جذر} لیں۔\\[0.3em]
   $\dfrac{\sqrt[3]{8}}{\sqrt[3]{27}}=\dfrac{2}{3}$}{\ealkey{Cube root} the \ealkey{numerator} and the \ealkey{denominator}.\\[0.3em]
   $\dfrac{\sqrt[3]{8}}{\sqrt[3]{27}}=\dfrac{2}{3}$}}

\qq[\Huge][\LARGE]{\ealpara{$0.04^{\frac{1}{2}}$ کا حساب لگائیں}{Work out $0.04^{\frac{1}{2}}$}}
  {\ealpara{$0.2\times 0.2=0.04$\\[0.3em]
   $0.04^{\frac{1}{2}}=0.2$}{$0.2\times 0.2=0.04$\\[0.3em]
   $0.04^{\frac{1}{2}}=0.2$}}

\qq[\Huge][\LARGE]{\ealpara{$\left(\dfrac{1}{32}\right)^{\frac{1}{5}}$ کا حساب لگائیں}{Work out $\left(\dfrac{1}{32}\right)^{\frac{1}{5}}$}}
  {\ealpara{$32$ کا \ealkeytr{پانچواں جذر} $2$ ہے، کیونکہ $2^{5}=32$۔\\[0.3em]
   $\dfrac{1}{2}$}{The \ealkey{fifth root} of $32$ is $2$, because $2^{5}=32$.\\[0.3em]
   $\dfrac{1}{2}$}}

%================================================================
\qtopic{\ealpara{$\frac{a}{b}$ کی شکل کے \ealkeytr{اس}}{\ealkey{Indices} of the form $\frac{a}{b}$}}
%================================================================

\tf[\LARGE][\large]{\ealpara{$8^{\frac{2}{3}}=4$}{$8^{\frac{2}{3}}=4$}}
  {\ealpara{\textbf{درست۔} \ealkeytr{اس} کا \ealkeytr{مخرج} \ealkeytr{جذر} دیتا ہے اور \ealkeytr{صورت} \ealkeytr{قوت} دیتا ہے۔
   \par\vspace{0.6em}
   نقشہ: \ealkeytr{مکعب جذر}، پھر \ealkeytr{مربع} کریں۔
   \par\vspace{0.6em}
   \flow{8}{\ealkey{cube root}}{2}{\ealkey{square} it}{4}}{\textbf{True.} The \ealkey{denominator} of the \ealkey{index} gives the \ealkey{root} and the \ealkey{numerator} gives the \ealkey{power}.
   \par\vspace{0.6em}
   Diagram: \ealkey{cube root}, then \ealkey{square} it.
   \par\vspace{0.6em}
   \flow{8}{\ealkey{cube root}}{2}{\ealkey{square} it}{4}}}

\tf[\LARGE][\large]{\ealpara{$16^{\frac{3}{4}}=12$}{$16^{\frac{3}{4}}=12$}}
  {\ealpara{\textbf{غلط۔} یہ $16$ کا $\frac{3}{4}$ ہے، $16$ کی \ealkeytr{قوت} $\frac{3}{4}$ نہیں۔
   \par\vspace{0.6em}
   نقشہ: \ealkeytr{چوتھا جذر}، پھر \ealkeytr{مکعب} کریں۔
   \par\vspace{0.6em}
   \flow{16}{\ealkey{fourth root}}{2}{\ealkey{cube} it}{8}}{\textbf{False.} This is $\frac{3}{4}$ of $16$, not $16$ to the \ealkey{power} $\frac{3}{4}$.
   \par\vspace{0.6em}
   Diagram: \ealkey{fourth root}, then \ealkey{cube} it.
   \par\vspace{0.6em}
   \flow{16}{\ealkey{fourth root}}{2}{\ealkey{cube} it}{8}}}

\tf{\ealpara{$25^{\frac{3}{2}}=75$}{$25^{\frac{3}{2}}=75$}}
  {\ealpara{\textbf{غلط۔} $\frac{3}{2}$ ایک \ealkeytr{اس} ہے، ضارب نہیں۔\\[0.3em]
   $\sqrt{25}=5$، پھر $5^{3}=125$۔}{\textbf{False.} The $\frac{3}{2}$ is an \ealkey{index}, not a multiplier.\\[0.3em]
   $\sqrt{25}=5$, then $5^{3}=125$.}}

\tf[\LARGE][\large]{\ealpara{$32^{\frac{3}{5}}$ نکالنے کے لیے آپ پہلے $\sqrt[5]{32}$ معلوم کر سکتے ہیں،\\[0.3em]
   پھر جواب کا \ealkeytr{مکعب} لے سکتے ہیں}{To work out $32^{\frac{3}{5}}$ you can find $\sqrt[5]{32}$ first,\\[0.3em]then \ealkey{cube} the answer}}
  {\ealpara{\textbf{درست۔} پہلے \ealkeytr{جذر} معلوم کرنے سے اعداد چھوٹے رہتے ہیں۔
   \par\vspace{0.6em}
   نقشہ: \ealkeytr{پانچواں جذر}، پھر \ealkeytr{مکعب} کریں۔
   \par\vspace{0.6em}
   \flow{32}{\ealkey{fifth root}}{2}{\ealkey{cube} it}{8}}{\textbf{True.} Finding the \ealkey{root} first keeps the numbers small.
   \par\vspace{0.6em}
   Diagram: \ealkey{fifth root}, then \ealkey{cube} it.
   \par\vspace{0.6em}
   \flow{32}{\ealkey{fifth root}}{2}{\ealkey{cube} it}{8}}}

\tf{\ealpara{$4^{\frac{3}{2}}$ اور $4^{1.5}$ کا مطلب ایک ہی ہے}{$4^{\frac{3}{2}}$ and $4^{1.5}$ mean the same thing}}
  {\ealpara{\textbf{درست۔} $\frac{3}{2}=1.5$، اس لیے دونوں $\sqrt{4}$ کا \ealkeytr{مکعب} ہیں، جو $8$ ہے۔}{\textbf{True.} $\frac{3}{2}=1.5$, so both are $\sqrt{4}$ \ealkey{cubed}, which is $8$.}}

\qq[\Huge][\large]{\ealpara{$4^{\frac{3}{2}}$ کا حساب لگائیں}{Work out $4^{\frac{3}{2}}$}}
  {\ealpara{نقشہ: \ealkeytr{جذر}، پھر \ealkeytr{مکعب} کریں۔
   \par\vspace{0.6em}
   \flow{4}{\ealkey{square root}}{2}{\ealkey{cube} it}{8}}{Diagram: \ealkey{square root}, then \ealkey{cube} it.
   \par\vspace{0.6em}
   \flow{4}{\ealkey{square root}}{2}{\ealkey{cube} it}{8}}}

\qq[\Huge][\large]{\ealpara{$8^{\frac{2}{3}}$ کا حساب لگائیں}{Work out $8^{\frac{2}{3}}$}}
  {\ealpara{نقشہ: \ealkeytr{مکعب جذر}، پھر \ealkeytr{مربع} کریں۔
   \par\vspace{0.6em}
   \flow{8}{\ealkey{cube root}}{2}{\ealkey{square} it}{4}}{Diagram: \ealkey{cube root}, then \ealkey{square} it.
   \par\vspace{0.6em}
   \flow{8}{\ealkey{cube root}}{2}{\ealkey{square} it}{4}}}

\qq[\Huge][\large]{\ealpara{$9^{\frac{3}{2}}$ کا حساب لگائیں}{Work out $9^{\frac{3}{2}}$}}
  {\ealpara{نقشہ: \ealkeytr{جذر}، پھر \ealkeytr{مکعب} کریں۔
   \par\vspace{0.6em}
   \flow{9}{\ealkey{square root}}{3}{\ealkey{cube} it}{27}}{Diagram: \ealkey{square root}, then \ealkey{cube} it.
   \par\vspace{0.6em}
   \flow{9}{\ealkey{square root}}{3}{\ealkey{cube} it}{27}}}

\qq[\Huge][\large]{\ealpara{$16^{\frac{3}{4}}$ کا حساب لگائیں}{Work out $16^{\frac{3}{4}}$}}
  {\ealpara{نقشہ: \ealkeytr{چوتھا جذر}، پھر \ealkeytr{مکعب} کریں۔
   \par\vspace{0.6em}
   \flow{16}{\ealkey{fourth root}}{2}{\ealkey{cube} it}{8}}{Diagram: \ealkey{fourth root}, then \ealkey{cube} it.
   \par\vspace{0.6em}
   \flow{16}{\ealkey{fourth root}}{2}{\ealkey{cube} it}{8}}}

\qq[\Huge][\large]{\ealpara{$27^{\frac{2}{3}}$ کا حساب لگائیں}{Work out $27^{\frac{2}{3}}$}}
  {\ealpara{نقشہ: \ealkeytr{مکعب جذر}، پھر \ealkeytr{مربع} کریں۔
   \par\vspace{0.6em}
   \flow{27}{\ealkey{cube root}}{3}{\ealkey{square} it}{9}}{Diagram: \ealkey{cube root}, then \ealkey{square} it.
   \par\vspace{0.6em}
   \flow{27}{\ealkey{cube root}}{3}{\ealkey{square} it}{9}}}

\qq[\Huge][\large]{\ealpara{$32^{\frac{3}{5}}$ کا حساب لگائیں}{Work out $32^{\frac{3}{5}}$}}
  {\ealpara{نقشہ: \ealkeytr{پانچواں جذر}، پھر \ealkeytr{مکعب} کریں۔
   \par\vspace{0.6em}
   \flow{32}{\ealkey{fifth root}}{2}{\ealkey{cube} it}{8}}{Diagram: \ealkey{fifth root}, then \ealkey{cube} it.
   \par\vspace{0.6em}
   \flow{32}{\ealkey{fifth root}}{2}{\ealkey{cube} it}{8}}}

\qq[\Huge][\large]{\ealpara{$125^{\frac{4}{3}}$ کا حساب لگائیں}{Work out $125^{\frac{4}{3}}$}}
  {\ealpara{نقشہ: \ealkeytr{مکعب جذر}، پھر چوتھی \ealkeytr{قوت}۔
   \par\vspace{0.6em}
   \flow{125}{\ealkey{cube root}}{5}{fourth \ealkey{power}}{625}}{Diagram: \ealkey{cube root}, then fourth \ealkey{power}.
   \par\vspace{0.6em}
   \flow{125}{\ealkey{cube root}}{5}{fourth \ealkey{power}}{625}}}

\qq[\Huge][\LARGE]{\ealpara{$\left(\dfrac{4}{9}\right)^{\frac{3}{2}}$ کا حساب لگائیں}{Work out $\left(\dfrac{4}{9}\right)^{\frac{3}{2}}$}}
  {\work{$\left(\dfrac{4}{9}\right)^{\frac{1}{2}}$ & $\dfrac{2}{3}$\\[0.6em]
         $\left(\dfrac{2}{3}\right)^{3}$ & $\dfrac{8}{27}$}}

\qq[\Huge][\LARGE]{\ealpara{$0.25^{\frac{3}{2}}$ کا حساب لگائیں}{Work out $0.25^{\frac{3}{2}}$}}
  {\ealpara{$\sqrt{0.25}=0.5$، کیونکہ $0.5^{2}=0.25$۔\\[0.3em]
   $0.5^{3}=0.125$}{$\sqrt{0.25}=0.5$, because $0.5^{2}=0.25$.\\[0.3em]
   $0.5^{3}=0.125$}}

\qq[\Huge][\LARGE]{\ealpara{$\left(\dfrac{16}{81}\right)^{\frac{3}{4}}$ کا حساب لگائیں}{Work out $\left(\dfrac{16}{81}\right)^{\frac{3}{4}}$}}
  {\work{$\left(\dfrac{16}{81}\right)^{\frac{1}{4}}$ & $\dfrac{2}{3}$\\[0.6em]
         $\left(\dfrac{2}{3}\right)^{3}$ & $\dfrac{8}{27}$}}

%================================================================
\qtopic{\ealpara{منفی \ealkeytr{کسری قوت نما} کے ساتھ \ealkeytr{اظہار} کو \ealkeytr{سادہ کرنا}}{\ealkey{Simplifying} \ealkey{expressions} with negative \ealkey{fractional index}es}}
%================================================================

\tf[\LARGE][\large]{\ealpara{$9^{-\frac{1}{2}}=\dfrac{1}{3}$}{$9^{-\frac{1}{2}}=\dfrac{1}{3}$}}
  {\ealpara{\textbf{درست۔} $\frac{1}{2}$ \ealkeytr{جذر} دیتا ہے اور منفی علامت \ealkeytr{مقلوب} دیتی ہے۔
   \par\vspace{0.6em}
   نقشہ: \ealkeytr{جذر}، پھر \ealkeytr{مقلوب}۔
   \par\vspace{0.6em}
   \flow{9}{\ealkey{square root}}{3}{\ealkey{reciprocal}}{\tfrac{1}{3}}}{\textbf{True.} The $\frac{1}{2}$ gives the \ealkey{square root} and the minus sign gives the \ealkey{reciprocal}.
   \par\vspace{0.6em}
   Diagram: \ealkey{square root}, then \ealkey{reciprocal}.
   \par\vspace{0.6em}
   \flow{9}{\ealkey{square root}}{3}{\ealkey{reciprocal}}{\tfrac{1}{3}}}}

\tf{\ealpara{$9^{-\frac{1}{2}}=-3$}{$9^{-\frac{1}{2}}=-3$}}
  {\ealpara{\textbf{غلط۔} منفی \ealkeytr{اس} \ealkeytr{مقلوب} دیتا ہے، منفی جواب نہیں۔\\[0.3em]
   $9^{-\frac{1}{2}}=\dfrac{1}{\sqrt{9}}=\dfrac{1}{3}$}{\textbf{False.} A negative \ealkey{index} gives a \ealkey{reciprocal}, not a negative answer.\\[0.3em]
   $9^{-\frac{1}{2}}=\dfrac{1}{\sqrt{9}}=\dfrac{1}{3}$}}

\tf[\LARGE][\large]{\ealpara{$8^{-\frac{2}{3}}=\dfrac{1}{4}$}{$8^{-\frac{2}{3}}=\dfrac{1}{4}$}}
  {\ealpara{\textbf{درست۔} \ealkeytr{مکعب جذر}، پھر \ealkeytr{مربع}، پھر \ealkeytr{مقلوب}۔
   \par\vspace{0.6em}
   نقشہ: \ealkeytr{مکعب جذر}، پھر \ealkeytr{مربع}، پھر \ealkeytr{مقلوب}۔
   \par\vspace{0.6em}
   \flowfour{8}{\ealkey{cube root}}{2}{\ealkey{square} it}{4}{\ealkey{reciprocal}}{\tfrac{1}{4}}}{\textbf{True.} \ealkey{Cube root}, then \ealkey{square}, then take the \ealkey{reciprocal}.
   \par\vspace{0.6em}
   Diagram: \ealkey{cube root}, then \ealkey{square}, then \ealkey{reciprocal}.
   \par\vspace{0.6em}
   \flowfour{8}{\ealkey{cube root}}{2}{\ealkey{square} it}{4}{\ealkey{reciprocal}}{\tfrac{1}{4}}}}

\tf{\ealpara{$16^{-\frac{3}{4}}=\dfrac{1}{12}$}{$16^{-\frac{3}{4}}=\dfrac{1}{12}$}}
  {\ealpara{\textbf{غلط۔} $\frac{3}{4}$ ایک \ealkeytr{اس} ہے، ضارب نہیں۔\\[0.3em]
   $\sqrt[4]{16}=2$، پھر $2^{3}=8$، اس لیے جواب $\dfrac{1}{8}$ ہے۔}{\textbf{False.} The $\frac{3}{4}$ is an \ealkey{index}, not a multiplier.\\[0.3em]
   $\sqrt[4]{16}=2$, then $2^{3}=8$, so the answer is $\dfrac{1}{8}$.}}

\tf[\LARGE][\large]{\ealpara{$25^{-\frac{1}{2}}$ $25^{-\frac{3}{2}}$ سے چھوٹا ہے}{$25^{-\frac{1}{2}}$ is smaller than $25^{-\frac{3}{2}}$}}
  {\ealpara{\textbf{غلط۔} \ealkeytr{اس} جتنا زیادہ منفی ہوگا، $1$ سے بڑے \ealkeytr{قاعدہ} کے لیے نتیجہ اتنا ہی چھوٹا ہوگا۔\\[0.3em]
   $25^{-\frac{1}{2}}=\dfrac{1}{5}=0.2$ اور $25^{-\frac{3}{2}}=\dfrac{1}{125}=0.008$۔}{\textbf{False.} The more negative the \ealkey{index}, the smaller the result for a \ealkey{base} bigger than $1$.\\[0.3em]
   $25^{-\frac{1}{2}}=\dfrac{1}{5}=0.2$ and $25^{-\frac{3}{2}}=\dfrac{1}{125}=0.008$.}}

\qq[\Huge][\LARGE]{\ealpara{$4^{-\frac{1}{2}}$ کا حساب لگائیں}{Work out $4^{-\frac{1}{2}}$}}
  {\ealpara{$4^{-\frac{1}{2}}=\dfrac{1}{\sqrt{4}}=\dfrac{1}{2}$}{$4^{-\frac{1}{2}}=\dfrac{1}{\sqrt{4}}=\dfrac{1}{2}$}}

\qq[\Huge][\LARGE]{\ealpara{$8^{-\frac{1}{3}}$ کا حساب لگائیں}{Work out $8^{-\frac{1}{3}}$}}
  {\ealpara{$8^{-\frac{1}{3}}=\dfrac{1}{\sqrt[3]{8}}=\dfrac{1}{2}$}{$8^{-\frac{1}{3}}=\dfrac{1}{\sqrt[3]{8}}=\dfrac{1}{2}$}}

\qq[\Huge][\LARGE]{\ealpara{$9^{-\frac{1}{2}}$ کا حساب لگائیں}{Work out $9^{-\frac{1}{2}}$}}
  {\ealpara{$9^{-\frac{1}{2}}=\dfrac{1}{\sqrt{9}}=\dfrac{1}{3}$}{$9^{-\frac{1}{2}}=\dfrac{1}{\sqrt{9}}=\dfrac{1}{3}$}}

\qq[\Huge][\large]{\ealpara{$27^{-\frac{2}{3}}$ کا حساب لگائیں}{Work out $27^{-\frac{2}{3}}$}}
  {\ealpara{نقشہ: \ealkeytr{مکعب جذر}، پھر \ealkeytr{مربع}، پھر \ealkeytr{مقلوب}۔
   \par\vspace{0.6em}
   \flowfour{27}{\ealkey{cube root}}{3}{\ealkey{square} it}{9}{\ealkey{reciprocal}}{\tfrac{1}{9}}}{Diagram: \ealkey{cube root}, then \ealkey{square}, then \ealkey{reciprocal}.
   \par\vspace{0.6em}
   \flowfour{27}{\ealkey{cube root}}{3}{\ealkey{square} it}{9}{\ealkey{reciprocal}}{\tfrac{1}{9}}}}

\qq[\Huge][\large]{\ealpara{$16^{-\frac{3}{4}}$ کا حساب لگائیں}{Work out $16^{-\frac{3}{4}}$}}
  {\ealpara{نقشہ: \ealkeytr{چوتھا جذر}، پھر \ealkeytr{مکعب}، پھر \ealkeytr{مقلوب}۔
   \par\vspace{0.6em}
   \flowfour{16}{\ealkey{fourth root}}{2}{\ealkey{cube} it}{8}{\ealkey{reciprocal}}{\tfrac{1}{8}}}{Diagram: \ealkey{fourth root}, then \ealkey{cube}, then \ealkey{reciprocal}.
   \par\vspace{0.6em}
   \flowfour{16}{\ealkey{fourth root}}{2}{\ealkey{cube} it}{8}{\ealkey{reciprocal}}{\tfrac{1}{8}}}}

\qq[\Huge][\LARGE]{\ealpara{$\left(\dfrac{1}{4}\right)^{-\frac{1}{2}}$ کا حساب لگائیں}{Work out $\left(\dfrac{1}{4}\right)^{-\frac{1}{2}}$}}
  {\ealpara{\ealkeytr{کسر} کا \ealkeytr{مقلوب} لیں، پھر \ealkeytr{جذر}۔\\[0.3em]
   $4^{\frac{1}{2}}=2$}{Take the \ealkey{reciprocal} of the \ealkey{fraction}, then \ealkey{square root}.\\[0.3em]
   $4^{\frac{1}{2}}=2$}}

\qq[\Huge][\LARGE]{\ealpara{$\left(\dfrac{9}{16}\right)^{-\frac{1}{2}}$ کا حساب لگائیں}{Work out $\left(\dfrac{9}{16}\right)^{-\frac{1}{2}}$}}
  {\ealpara{\ealkeytr{کسر} کا \ealkeytr{مقلوب} لیں، پھر \ealkeytr{جذر}۔\\[0.4em]
   $\left(\dfrac{16}{9}\right)^{\frac{1}{2}}=\dfrac{4}{3}$}{Take the \ealkey{reciprocal} of the \ealkey{fraction}, then \ealkey{square root}.\\[0.4em]
   $\left(\dfrac{16}{9}\right)^{\frac{1}{2}}=\dfrac{4}{3}$}}

\qq[\Huge][\large]{\ealpara{$\left(\dfrac{8}{27}\right)^{-\frac{2}{3}}$ کا حساب لگائیں}{Work out $\left(\dfrac{8}{27}\right)^{-\frac{2}{3}}$}}
  {\work{$\left(\dfrac{8}{27}\right)^{-\frac{2}{3}}$ & $\left(\dfrac{27}{8}\right)^{\frac{2}{3}}$\\[0.6em]
         $\left(\dfrac{27}{8}\right)^{\frac{2}{3}}$ & $\left(\dfrac{3}{2}\right)^{2}$\\[0.6em]
         $\left(\dfrac{3}{2}\right)^{2}$ & $\dfrac{9}{4}$}}

\qq[\Huge][\LARGE]{\ealpara{\ealkeytr{سادہ کرنا} $\dfrac{x^{\frac{1}{2}}}{x^{-\frac{3}{2}}}$}{\ealkey{Simplify} $\dfrac{x^{\frac{1}{2}}}{x^{-\frac{3}{2}}}$}}
  {\ealpara{\ealkeytr{اس} تفریق ہوتے ہیں۔\\[0.3em]
   $x^{\frac{1}{2}-\left(-\frac{3}{2}\right)}=x^{2}$}{The \ealkey{indices} are subtracted.\\[0.3em]
   $x^{\frac{1}{2}-\left(-\frac{3}{2}\right)}=x^{2}$}}

\qq[\LARGE][\LARGE]{\ealpara{\ealkeytr{سادہ کرنا} $\left(16x^{8}\right)^{-\frac{3}{4}}$}{\ealkey{Simplify} $\left(16x^{8}\right)^{-\frac{3}{4}}$}}
  {\work{$\left(16x^{8}\right)^{\frac{1}{4}}$ & $2x^{2}$\\[0.4em]
         $\left(2x^{2}\right)^{3}$ & $8x^{6}$\\[0.4em]
         $\left(16x^{8}\right)^{-\frac{3}{4}}$ & $\dfrac{1}{8x^{6}}$}}

\end{document}